\documentclass[12pt,a4paper]{article}
\usepackage[utf8]{inputenc}
\usepackage[T2A]{fontenc}
\usepackage[russian]{babel}
\usepackage{amsmath, amssymb, amsthm}
\usepackage{tikz}
\usepackage{pgfplots}
\pgfplotsset{compat=1.18}

\title{Экспонента: определение свойства и применение}
\author{insti2t}
\date{\today}

\begin{document}

\maketitle

\section{Определение экспоненты}
Экспоненциальная функция (экспонента) определяется как
\begin{equation} \label{1}
	e^x=exp(x)
\end{equation}
где 
\begin{equation}\label{2}
 e\approx2,71828
\end{equation}
число Эйлера, основание натурального логорифма.
\\
Альтернативное определение - через предел:
\begin{equation}\label{3}
	 e^x=\lim_{n \to \infty}(1+\frac{x}{n})^n
\end{equation}

\section{Разложение в ряд Тейлора}
Одно из важнейших представлений экспоненты - степенной ряд:
\begin{equation}\label{4}
	e^x=\sum_{n=0}^{\infty}\frac{x^n}{n!}=1+x+\frac{x^2}{2!}+\frac{x^3}{3!}+...
\end{equation}
Этот ряд сходится при всех $x \in \mathbb{R}$
Таким образом эти формулы крутые - \ref{1} \ref{2} \ref{3} \ref{4}
\section{Основные свойства}
Перечислим ключевые свойства экспоненты:
\begin{itemize}
	\item Мультипликативность: $e^{x+y}=e^x*e^y$
	\item Производная равна самой функции: $\frac{d}{dx}e^x=e^x$
	\item Интеграл: $\int e^xdx=e^x+C$
	\\
	\\
	\item Обратная функция - натуральный логарифм: $ln(e^x)=x$
	\item Экспонента всегда положительна: $e^x > 0$ для всех $x \in \mathbb{R}$
\end{itemize}
\section{Дифференциальное уравнение}
Экспонента является решением задачи Коши:
\begin{equation}\label{5}
	\frac{dy}{dx}=y, y(0)=1
\end{equation}
Решение: $y(x)=e^x$. Это свойство лежит в основе описания процессов экспоненциального роста и распада.
\section{График экспоненциальной функции}
Ниже - код для построения графика $y=e^x$ средствами TikZ/pfgplots:
\begin{figure}[h!]
\centering
\begin{tikzpicture}
\begin{axis}[
width=0.7\textwidth,
height=0.5\textwidth,
axis lines=middle,
xlabel={$x$},
ylabel={$y$},
xmin=-3, xmax=3,
ymin=-1, ymax=21,
grid=both,
samples=150,
domain=-3:3,
legend pos=north west
]

\addplot[blue, thick]{exp(x)};
\addlegendentry{$y=e^x$}
\end{axis}
\end{tikzpicture}
\caption{График функций $y=e^x$}
\label{fig:exp}
\end{figure}
Таким образом с помощью формулы \ref{5} мы сделали чтото вроде график
\section{Примеры применения}
\subsection{Непрерывный рост и распад}
Модель экспоненциального роста:
$N(t)=N_{0}e^{kt}$
где:
\begin{itemize}
\item$N_{0}$-начальное значение,
\item$k$-коэффицент роста $(k>0)$ или распада $(k<0)$,
\item$t$-время.
Примеры: рост популяции, радиоактивный распад, сложные проценты.
\end{itemize}
\subsection{Комплексная экспонента и формула Эйлера}
Для комплексных чисел справедлива формула Эйлера:
\begin{equation}\label{6}
e^{ix}=\cos x+i\sin x
\end{equation}
Отсюда следует знаменитая тождественная формула:
\begin{equation}\label{7}
e^{i\pi}+1=0
\end{equation}
Таким образом формулы  \ref{6} \ref{7} тоже крутые!
\section{Заключение}
Экспонента - одна из важнейших функций в математике и ее приложениях. Ее уникальные свойства (совпадение функций и производной, мультипликативность, связь, с комплексными числами) делают ее незаменимой в анализе, физике, экономике и инженерии.  

\end{document}
